\section{Introduction}
Coverage design has three logically different tasks. One must establish that a specified target is covered, determine which failures destroy that property, and choose a design or protection policy that meets a quantitative requirement. Computing the topology of the union addresses an important part of the geometry, but it does not by itself answer the target question. A contractible covered region can be disjoint from the asset that must be monitored.

This paper studies a finite model in which the geometry can be certified exactly. Physical footprints are closed disks, the target is a finite union of closed nondegenerate triangles, and failure rules are monotone with respect to additional functioning contacts. We use the model to connect a canonical geometric certificate with exact failure budgets, common-radius design, uncertainty and fault-domain assignment. Every statement below specifies its hypotheses; no physical noise distribution or hardware cost law is inferred from the geometry.

\subsection{Two reusable certificates}
For a point $x$ of the target $T$, let $S(x)$ be the set of contacts whose footprints contain $x$. Define
\begin{equation}\label{eq:overview-antichain}
\mathcal A_T=\Min\{S(x):x\in T\},
\end{equation}
where $\Min$ removes every strict superset of another member. For the regular closed targets considered here, every member has a representative away from all disk and target boundaries. A vertical-slice construction finds a rational representative for each member without approximate intersection decisions.

The first certificate is this antichain together with its witnesses. If $\F$ is any downward-closed family of allowed failures, then
\begin{equation}\label{eq:overview-universal}
T\subseteq R_{\F}\quad\Longleftrightarrow\quad
A\notin\F\text{ for every }A\in\mathcal A_T.
\end{equation}
One geometric computation therefore answers every such failure query. The antichain is the minimal-cut family of the target-coverage reliability system, rather than a new definition of a coherent system. Its geometric origin makes it possible to attach an explicit target location to each cut.

The second certificate concerns common-radius design. Fix sites and target but let all primitive disk radii equal an unknown $r$. We construct a finite set $Q_T$ of rational points, independent of the winning rule. For every monotone rule $w$ that wins on the full contact set, define
\[
\rho_w(x)=\min_{\sigma:\,w(\sigma)=1}\max_{i\in\sigma}\|x-c_i\|.
\]
Then the exact critical radius satisfies
\begin{equation}\label{eq:overview-radius}
(r_w^*)^2=\max_{q\in Q_T}\rho_w(q)^2\in\mathbb Q.
\end{equation}
The smaller support construction uses target vertices, bisector--edge intersections and three-site circumcenters, for $O(n^3+n^2h+h)$ candidates. It is valid simultaneously for all rules and all common radii. A fixed-dimensional theorem uses at most $d+1$ supporting sites on a $d$-dimensional target. We retain the larger full distance-order arrangement as an independent comparison; no optimal running-time claim is made.

\subsection{Pure mathematics, applied mathematics and engineering}
The pure mathematical part identifies the canonical antichain, proves its minimality for universal failure queries, describes its behavior under monotone label maps and relates the geometry to nerve and Dowker models. The general support theorem covers functions with a common continuous convex part and affine differences. Its finite support set works for every monotone rule and every common convex term. The radius function is a lattice polynomial in the contact distances, an established class of functions~\cite{lattice,reliability-lattice}.

The applied mathematical part gives exact arithmetic, finite algorithms, geometric uncertainty reductions and optimization certificates. For product uncertainty sets, robustification commutes with every monotone Boolean contact rule. Independent displacement balls therefore reduce exactly to disk erosion. Independent polygonal displacement sets reduce to finite intersections of disks at their vertices. Protection allocation is a linear program whose solutions are checked with rational primal and dual witnesses.

The engineering part derives complete target coverage, failure budgets, wiring choices and uncertain placement in exact symbolic examples. In the four-corner example, changing the domain assignment changes the model's coverage reliability from $81/100$ to $99/100$ at domain failure probability $1/10$, while leaving positions, radii and contact count unchanged. The exact squared radii for surviving a domain loss change from $89/25$ to $68/25$. These are analytical consequences of the stated model, not experimental performance claims about manufactured sensors.

\begin{figure}[tb]
\centering
\begin{tikzpicture}[>=Latex,b/.style={draw=blue!50!black,rounded corners,align=center,font=\small,inner sep=5pt},node distance=9mm and 7mm]
\node[b] (geo) {Sites and target\\exact geometry};
\node[b,right=of geo] (cuts) {Target antichain\\weakest-point witnesses};
\node[b,right=of cuts] (radius) {Distance-order vertices\\exact critical radii};
\node[b,below=of geo] (logic) {Pure mathematics\\monotone logic and topology};
\node[b,below=of cuts] (opt) {Applied mathematics\\uncertainty and duality};
\node[b,below=of radius] (eng) {Engineering\\budgets and wiring};
\draw[->] (geo)--(cuts);\draw[->] (geo) to[bend left=20] (radius);
\draw[->] (cuts)--(logic);\draw[->] (cuts)--(opt);\draw[->] (radius)--(eng);\draw[->] (opt)--(eng);
\end{tikzpicture}
\caption{The two finite geometric certificates and their roles. Geometry is evaluated before selecting a failure rule or protection allocation.}
\end{figure}

\subsection{Relation to established work}
Arrangements of algebraic curves and exact geometric predicates are standard tools in computational geometry~\cite{arrangements}. Coverage algorithms based on circle intersections and local geometry precede this paper~\cite{integrated,coverage2003}. Homological sensor criteria and multicover nerves are likewise established~\cite{ghrist,desilva,controlled,sheehy,corbet,alonso}. We give the needed maps and proofs because the failure semantics and boundary conventions matter for the certificates.

Minimal cut sets and structure functions belong to reliability theory~\cite{barlow}. Cover decomposition and polychromatic hypergraph coloring provide the appropriate context for fault-domain assignment~\cite{polychromatic}. Max--min scalar functions built from monotone logic are lattice polynomials~\cite{lattice,reliability-lattice}; Boolean combinations of balls and generalized Voronoi decompositions also occur in the geometric work of Csik\'os~\cite{csikos}. The protection program uses ordinary linear programming duality. The treatment of uncertain geometry is situated within robust optimization~\cite{robustbook}, and marginal-only probability bounds are instances of the logical probability programs studied by Hailperin~\cite{hailperin}.

Accordingly, the contribution proposed here is the precise geometric construction and its verified synthesis, including reusable rational witnesses, exact uncertainty reductions and common-radius certificates under one failure interface. The manuscript does not claim to have discovered arrangements, minimal cuts, lattice polynomials or LP duality. Establishing priority for the full combination and assessing its significance require specialist comparison with this literature.

The earlier sensor manuscript~\cite{original} provides background for the contact-intersection model. The ECoG archive~\cite{ecog} motivates separating exact predicates, continuous target certificates and illustrations. The authors' circuit-topos, topological-deformation and linkage studies~\cite{topoi,suspension,linkage} are contextual references; none is a premise of a proof here.

\subsection{Conventions and scope}
All targets in the exact planar algorithms are finite unions of nondegenerate closed triangles with rational vertices. These targets are regular closed: $T=\overline{\operatorname{int}T}$. They may be disconnected or have holes. A wire target or an isolated point requires a lower-dimensional treatment and is excluded from the full-dimensional sampling theorem.\footnote{Regular closed does not mean convex or simply connected. The condition ensures that failure of a closed coverage set at a target boundary point is witnessed arbitrarily nearby in the target interior.}

Disk centers and radii are rational. A supplied binary64 number is treated as its exact rational value. The new target and radius routines preserve explicit \texttt{Fraction} inputs; the legacy topology routine documents its separate binary64 conversion convention. All masks retain physical contact labels, including coincident footprints. Ordinary homology is over $\FF$ unless another field is specified. The target certificates themselves do not depend on a choice of homology coefficients.

The main text develops the exact certificates and their uses. The appendices include comparison homotopies, the closed-cover chain argument, arithmetic derivations, dual certificates, worked examples and boundary counterexamples. Every mathematical conclusion in the manuscript is proved analytically. Optional software checks are supplied separately and are not premises of those proofs. Figures illustrate the arguments; no grid resolution or simulated observation determines a theorem or example.
